%--------------------------------------------------------------------------
%	PACKAGES AND OTHER DOCUMENT CONFIGURATIONS
%--------------------------------------------------------------------------
\documentclass[12pt,a4paper]{article}
\usepackage[utf8]{inputenc}
\usepackage[english]{babel}
\usepackage[T1]{fontenc}
\usepackage{amsmath}
\usepackage{amsfonts}
\usepackage{amssymb}
\usepackage{graphicx}
\usepackage{lmodern}
\usepackage[left=3cm,right=3cm,top=2.5cm,bottom=2.5cm]{geometry}

\usepackage{fancyhdr} % Required for custom headers
\usepackage{lastpage} % Required to determine the last page for the footer
\usepackage{extramarks} % Required for headers and footers
\usepackage[usenames,dvipsnames]{color} % Required for custom colors
\usepackage{graphicx} % Required to insert images
\usepackage{caption}
\usepackage{subcaption}
\usepackage{listings} % Required for insertion of code
%\usepackage{courier} % Required for the courier font
\usepackage{verbatim}
\usepackage{multirow}
\usepackage{eurosym}
\usepackage{url}
\usepackage{hyperref}

\usepackage{todonotes}

\lstset{
	language=bash,
	tabsize=4,
	rulecolor=,
    basicstyle=\scriptsize,
    upquote=true,
    aboveskip={1.5\baselineskip},
    columns=fixed,
    showstringspaces=false,
    extendedchars=true,
    breaklines=true,
    %prebreak = \raisebox{0ex}[0ex][0ex]{\ensuremath{\hookleftarrow}},
    frame=single,
    showtabs=false,
    showspaces=false,
    showstringspaces=false,
    numbers=left,
    stepnumber=0
}

\setlength\parindent{0pt} % Removes all indentation from paragraphs

% sections in Alph
\renewcommand{\thesection}{\Roman{section}}

\begin{document}
	
%--------------------------------------------------------------------------
%	TITLE PAGE
%--------------------------------------------------------------------------
\begin{titlepage}
\newcommand{\HRule}{\rule{\linewidth}{0.5mm}} % Defines a new command for the horizontal lines, change thickness here
\centering % Center everything on the page
 
%	HEADING SECTIONS
\begin{figure}[!h]
	\begin{center}
	\includegraphics[height=4cm]{liu.jpg}
	\end{center}
\end{figure}

%\null
%\vspace{1cm}
%\textsc{\Large Linköping University}\\[1cm] % Name of your university/college
\textsc{\large Department of Computer and Information Science}\\[0.5cm] % Major heading such as course name
\textsc{\large TDTS06 Computer Networks}\\[2.5cm] % Minor heading such as course title

%	TITLE SECTION

\HRule \\[0.4cm]
{ \LARGE \bfseries Distance Vector Routing\\[0.4cm] % Title of your document
\Large \bfseries Report} \\[0.4cm]

\HRule \\[1.5cm]


%	AUTHOR SECTION

\large
\begin{centering}
Group C/D 13 \\[0.2cm]
\end{centering}
{\begin{tabular}{lll}
\textsc{Chvátal} & Martin & march011@student.liu.se\\
\textsc{Peschke} & Lena & lenpe782@student.liu.se\\
\end{tabular}}
\\[1.5cm]

\normalsize
Teaching assistant :\\[0.2cm]
{\begin{tabular}{ll}
\textsc{Schmidt} Johannes & johannes.schmidt@liu.se \\
\end{tabular}}
\\[2cm]

%	DATE SECTION
\vfill
{\normalsize \today} % Date, change the \today to a set date if you want to be precise

\newpage

\end{titlepage}

%--------------------------------------------------------------------------
%	TABLE OF CONTENTS
%--------------------------------------------------------------------------

%\pagenumbering{gobble}
%\clearpage
%\thispagestyle{empty}
%\tableofcontents
%\clearpage
%\pagenumbering{arabic}

%--------------------------------------------------------------------------
%	CONTENT
%--------------------------------------------------------------------------
  
% i
\section{How distance vector routing works}
\paragraph{In general}
Distance vector routing uses a decentralized routing algorithm to find the least-cost path between routers in a network.
The distance vector routing algorithm considers a network as an undirected graph of routers, represented as nodes, and links, represented as edges. It is iterative, asynchronous and distributed \cite[p.~371]{cn}. It makes use of the Bellman-Ford equation:
\begin{equation}
d_x(y) = \min_v\{c(x,v) + d_v(y)\}
\label{eq:bf}
\end{equation}
The equation states that the cost of the least-cost path from node $x$ to $y$, $d_x(y)$, is the minimum of the sum of the cost to get from $x$ to a neighbor $v$ and the cost of the least-cost path from $v$ to $y$.

The algorithm works by making the nodes exchange distance vectors.
To make use of the equation, each node needs to know about all of the other nodes in the network and the cost to its direct neighbors. Let's walk through the algorithm step by step.

\begin{enumerate}
\item Initially, each node knows the cost to its direct neighbors. It sends its distance vector to its neighbors.

\item Upon receiving an update (link cost change or new distance vector from a neighbor), a node recomputes its distance vector using the Bellman-Ford equation~(\ref{eq:bf}). If there is a change, it sends its own new vector to the all the neighbors.

\item ``The process of receiving updated distance vectors from neighbors, recomputing routing table entries, and informing neighbors of changed costs of the least-cost path to a destination continues until no update messages are sent.''~\cite[p.~375]{cn}
\end{enumerate}

After step 3, the estimated costs have converged to the actual least-costs.

% ii
\section{How we tested the algorithm}
\todo[inline]{wait for the code}

% iii
\section{Poisoned reverse}
% explain what it is
When a link cost in the network increases, a count-to-infinity problem can occur. A packet can then be caught in a routing loop until the nodes have figured out the shortest path, which can take a while (as many iterations as the difference between the previous and the new least-cost). Poisoned reverse is an addition to the DVR algorithm which helps avoiding this situation. If a node $x$ routes through $y$ in order to get to $z$, $x$ lies to $y$ about its distance to $z$, so that an increase in the cost link between $y$ and $z$ results in an convergent update that takes only two iterations. 

\paragraph{Where it fails}
% example topology
Poisonous reverse fails if a count-to-infinity problem triggers a loop of three or more nodes \cite[p.~378]{cn}.

\paragraph{Solutions}
\todo[inline]{find solutions}
% find them

\bibliographystyle{alpha}
\bibliography{biblio}
\nocite{*}
    
\end{document}